% ============================================================
%  Преамбула учебника «Основы Docker»
%  Компиляция: lualatex (см. Makefile)
% ============================================================

% ---------- Шрифты и язык ----------
\usepackage{fontspec}
\setmainfont{PT Serif}[Ligatures=TeX]
\setsansfont{PT Sans}[Ligatures=TeX]
\setmonofont{PT Mono}[Scale=0.90, Ligatures=TeXOff]

\usepackage[english,russian]{babel}
\usepackage[autostyle]{csquotes}

% ---------- Геометрия страницы ----------
\usepackage[
  a4paper,
  top=2.5cm, bottom=2.7cm,
  left=2.7cm, right=2.3cm,
  headsep=0.8cm, footskip=1.2cm
]{geometry}
\raggedbottom
\setlength{\parindent}{1.25em}
\setlength{\parskip}{0.25em}
\linespread{1.04}
\emergencystretch=3em

% ---------- Математика и таблицы ----------
\usepackage{amsmath}
\usepackage{amssymb}
\usepackage{booktabs}
\usepackage{tabularx}
\usepackage{array}
\usepackage{longtable}
\usepackage{multirow}
\usepackage{enumitem}
\setlist{itemsep=0.15em, topsep=0.35em, parsep=0.2em}
\renewcommand{\arraystretch}{1.18}

% ---------- Графика ----------
\usepackage{graphicx}
\usepackage{xcolor}
\usepackage{tikz}
\usetikzlibrary{
  arrows.meta, positioning, shapes.geometric, shapes.misc,
  fit, backgrounds, calc, chains, shadows.blur,
  decorations.pathreplacing
}

% ---------- Палитра ----------
\definecolor{ink}{HTML}{1F2328}
\definecolor{accent}{HTML}{14507A}      % основной синий
\definecolor{accent2}{HTML}{0E86B8}     % голубой акцент
\definecolor{accentpale}{HTML}{EAF3F9}
\definecolor{warm}{HTML}{B45309}        % оранжевый
\definecolor{good}{HTML}{1B6E4A}        % зелёный
\definecolor{violet}{HTML}{5A3E9B}      % фиолетовый
\definecolor{muted}{HTML}{6A737D}
\definecolor{codebg}{HTML}{F6F8FA}
\definecolor{codeframe}{HTML}{D8DEE4}
\definecolor{rulegray}{HTML}{D0D7DE}
\color{ink}

% ---------- Гиперссылки ----------
\usepackage[hidelinks]{hyperref}
\hypersetup{
  colorlinks=true,
  linkcolor=accent,
  urlcolor=accent2,
  citecolor=accent,
  pdftitle={Основы Docker. Учебник},
  pdfsubject={Контейнеризация, Docker Engine, Docker Compose},
  pdfkeywords={Docker, контейнеры, образы, Dockerfile, Compose, DevOps},
  bookmarksnumbered=true,
  bookmarksopen=true
}
\usepackage{bookmark}

% ---------- Микротипографика ----------
\usepackage{microtype}

% ---------- Заголовки ----------
\usepackage{titlesec}
\titleformat{\chapter}[display]
  {\sffamily\bfseries\raggedright}
  {\normalsize\color{accent2}\MakeUppercase{\chaptertitlename}\ \thechapter}
  {6pt}
  {\fontsize{25}{29}\selectfont\color{accent}}
  [\vspace{8pt}{\color{accent2}\rule{\textwidth}{1.1pt}}]
\titlespacing*{\chapter}{0pt}{-6pt}{26pt}

\titleformat{\section}
  {\sffamily\bfseries\large\color{accent}}
  {\thesection}{0.6em}{}
\titlespacing*{\section}{0pt}{1.6em}{0.55em}

\titleformat{\subsection}
  {\sffamily\bfseries\normalsize\color{accent!90!black}}
  {\thesubsection}{0.5em}{}
\titlespacing*{\subsection}{0pt}{1.2em}{0.4em}

\titleformat{\subsubsection}[runin]
  {\sffamily\bfseries\normalsize\color{ink}}
  {\thesubsubsection}{0.5em}{}[.\enspace]

\titleformat{\part}[display]
  {\sffamily\bfseries\centering}
  {\Large\color{accent2}\MakeUppercase{\partname}\ \thepart}
  {14pt}
  {\fontsize{30}{34}\selectfont\color{accent}}
  [\vspace{10pt}{\color{accent2}\rule{0.5\textwidth}{1.1pt}}]

% ---------- Колонтитулы ----------
\usepackage{fancyhdr}
\pagestyle{fancy}
\fancyhf{}
\fancyhead[L]{\small\sffamily\color{muted}\nouppercase{\leftmark}}
\fancyhead[R]{\small\sffamily\color{muted}\thepage}
\setlength{\headheight}{22pt}
\renewcommand{\headrulewidth}{0.4pt}
\renewcommand{\headrule}{\hbox to\headwidth{\color{rulegray}\leaders\hrule height \headrulewidth\hfill}}
\fancypagestyle{plain}{%
  \fancyhf{}%
  \fancyfoot[C]{\small\sffamily\color{muted}\thepage}%
  \renewcommand{\headrulewidth}{0pt}%
}
\renewcommand{\chaptermark}[1]{\markboth{\thechapter.\ #1}{}}
\renewcommand{\sectionmark}[1]{}

% ---------- Подписи ----------
\usepackage[font=small,labelfont=bf,labelsep=period]{caption}

% ---------- Коробки ----------
\usepackage[most]{tcolorbox}
\tcbset{
  boxbase/.style={
    enhanced,
    breakable,
    boxrule=0.4pt,
    arc=2pt,
    left=10pt, right=10pt, top=7pt, bottom=7pt,
    before skip=10pt, after skip=10pt,
    fonttitle=\sffamily\bfseries,
    coltitle=accent!75!black,
    attach title to upper={\par\vspace{0.25em}},
  },
  codebase/.style={
    enhanced,
    breakable,
    boxrule=0.4pt,
    arc=2pt,
    left=8pt, right=8pt, top=6pt, bottom=6pt,
    before skip=9pt, after skip=9pt,
  }
}

\newtcolorbox{defbox}[1][]{%
  boxbase,
  colback=accent!5,
  colframe=accent!45!black,
  colbacktitle=accent!85!black,
  title={\sffamily\bfseries Определение},
  #1
}

\newtcolorbox{notebox}[1][]{%
  boxbase,
  colback=black!3,
  colframe=black!35,
  colbacktitle=black!55,
  title={\sffamily\bfseries Замечание},
  #1
}

\newtcolorbox{warnbox}[1][]{%
  boxbase,
  colback=warm!6,
  colframe=warm!70!black,
  colbacktitle=warm!85!black,
  title={\sffamily\bfseries Важно},
  #1
}

\newtcolorbox{tipbox}[1][]{%
  boxbase,
  colback=good!6,
  colframe=good!65!black,
  colbacktitle=good!80!black,
  title={\sffamily\bfseries Совет},
  #1
}

\newtcolorbox[auto counter, number within=chapter]{labbox}[2][]{%
  boxbase,
  colback=violet!5,
  colframe=violet!65!black,
  colbacktitle=violet!85!black,
  title={\sffamily\bfseries Лабораторная работа \thetcbcounter. #2},
  #1
}

\newtcolorbox{summarybox}[1][]{%
  boxbase,
  colback=accent2!7,
  colframe=accent2!70!black,
  colbacktitle=accent2!85!black,
  title={\sffamily\bfseries Итоги главы},
  #1
}

% ---------- Листинги ----------
\usepackage{listings}
\lstdefinestyle{dev}{
  basicstyle=\ttfamily\footnotesize,
  keywordstyle=\color{accent}\bfseries,
  commentstyle=\color{good!55!black},
  stringstyle=\color{warm},
  numbers=none,
  breaklines=true,
  breakatwhitespace=false,
  breakindent=1.5em,
  postbreak=\mbox{\textcolor{accent2}{$\hookrightarrow$}\space},
  columns=fullflexible,
  keepspaces=true,
  showstringspaces=false,
  tabsize=2,
  upquote=true,
  extendedchars=true,
  xleftmargin=0pt,
  aboveskip=0pt,
  belowskip=0pt,
}

\lstdefinelanguage{dockerfile}{
  sensitive=true,
  morekeywords={FROM,RUN,CMD,LABEL,MAINTAINER,EXPOSE,ENV,ADD,COPY,ENTRYPOINT,VOLUME,USER,WORKDIR,ARG,ONBUILD,STOPSIGNAL,HEALTHCHECK,SHELL,AS},
  morecomment=[l]{\#},
}

\lstdefinelanguage{composeyaml}{
  sensitive=false,
  morekeywords={services,volumes,networks,image,build,container_name,ports,environment,env_file,depends_on,restart,healthcheck,command,entrypoint,profiles,secrets,configs,deploy,resources,limits,reservations,watch,develop,condition,test,interval,timeout,retries,start_period,target,context,dockerfile,args},
  morecomment=[l]{\#},
  morestring=[b]",
  morestring=[b]',
}

\lstdefinelanguage{plain}{}

\newcommand{\codesetup}{%
  \lstset{style=dev,%
    backgroundcolor=\color{codebg},%
    frame=leftline,%
    framerule=2.2pt,%
    rulecolor=\color{accent2!75!black},%
    framesep=7pt,%
    xleftmargin=11pt,%
    xrightmargin=4pt,%
    aboveskip=9pt,%
    belowskip=9pt,%
  }%
}

\lstnewenvironment{code}[1][]{\codesetup\lstset{#1}}{}
\lstnewenvironment{shell}[1][]{\codesetup\lstset{#1}}{}

% ---------- Команды ----------
\newcommand{\cmd}[1]{\texttt{#1}}
\newcommand{\file}[1]{\texttt{#1}}
\newcommand{\dcmd}[1]{\texttt{docker #1}}
\newcommand{\envvar}[1]{\texttt{#1}}
\newcommand{\flag}[1]{\texttt{#1}}

% Ненумерованная глава с записью в содержание и колонтитул
\newcommand{\frontchapter}[1]{%
  \chapter*{#1}%
  \addcontentsline{toc}{chapter}{#1}%
  \markboth{#1}{}%
}

% Вопросы для самопроверки
\newenvironment{selfcheck}
  {\begin{enumerate}[label=\arabic*., leftmargin=2em, itemsep=0.3em]}
  {\end{enumerate}}

\setlength{\overfullrule}{6pt}
